\section{Wo die Type-I-Codebooks erzeugt werden}

Die Type-I-Codebooks sind in zwei Schichten organisiert. Die Matrix-Mathematik liegt in \file{lib/ran/precoding/},
die Erzeugung des kompletten Codebook-Repositories für eine Zelle im FAPI-Adaptor.

\subsection{Die Mathematik: eine Matrix aus einem PMI}

\begin{itemize}
  \item \file{lib/ran/precoding/type1\_sp/precoding\_codebooks.cpp} ist der Kern. \code{make\_type1\_sp\_mode1()}
    baut die Gewichtsmatrix nach TS\,38.214 Tabellen 5.2.2.2.1-5 bis -8 für 1 bis 8 Layer. Pro Layer wird
    \code{make\_layer\_type1\_sp\_mode1()} aufgerufen: Kronecker-Produkt der horizontalen und vertikalen
    DFT-Vektoren, die zweite Polarisation erhält den Co-Phasing-Faktor $\varphi$.
    \code{get\_k1\_and\_k2\_2\_layers()} und \code{get\_k1\_and\_k2\_3\_4\_layers()} sind die Tabellen
    5.2.2.2.1-3 und -4 (Beam-Offset zwischen Layern über $i_{1,3}$).
  \item Dieselbe Datei enthält \code{calculate\_mimo\_matrix()} und \code{get\_beams\_from\_pmi()}, die zerlegte
    Variante für Topologien mit Beam-Raster (Abschnitt~\ref{sec:zwei-matrizen}).
  \item \file{lib/ran/precoding/two\_port/precoding\_codebooks.cpp}: Sonderfall 2 Ports, Tabelle 5.2.2.2.1-1,
    hart codiert.
  \item \file{lib/ran/precoding/precoding\_codebooks.cpp}: Dispatcher. \code{make\_precoding()} und
    \code{get\_mimo\_matrix\_from\_pmi()} machen ein \code{std::visit} über die PMI-Variante
    (Abschnitt~\ref{sec:dispatcher}).
  \item \file{lib/ran/precoding/type1\_sp/precoding\_codebook\_helpers.cpp}: Wertebereiche und Bitbreiten der
    PMI-Felder $i_{1,1}, i_{1,2}, i_{1,3}, i_2$ je Panel und Rang, TS\,38.212 Tabelle 6.3.1.1.2-1.
  \item \file{lib/ran/precoding/precoding\_codebook\_configuration.cpp}: \code{get\_single\_panel\_info()}, die
    Tabelle $N_1, N_2, O_1, O_2$ der unterstützten Panels $2\times1$, $2\times2$, $4\times1$.
\end{itemize}

Die Typen stehen in \file{include/ocudu/ran/precoding/}: \file{precoding\_matrix\_indicator.h} mit
\code{pmi\_typeI\_single\_panel}, \file{precoding\_codebook\_configuration.h} und
\file{precoding\_codebook\_type1\_helpers.h}.

\begin{lstlisting}[language=C++,title={include/ocudu/ran/precoding/precoding\_matrix\_indicator.h}]
struct pmi_typeI_single_panel {
  /// Single-panel topology configuration.
  pmi_codebook_typeI_single_panel panel_config;
  /// PMI parameter i_{1,1}. Values: {0,...,N_1*O_1 - 1}, at most {0,...,63}.
  uint8_t i_1_1;
  /// PMI parameter i_{1,2}. Only available for N_2 > 1 or v > 2.
  std::optional<uint8_t> i_1_2;
  /// PMI parameter i_{1,3}. Only available for v in {2,3,4}.
  std::optional<uint8_t> i_1_3;
  /// PMI parameter i_2.
  uint8_t i_2;
};

using precoding_matrix_indicator =
    std::variant<std::monostate, pmi_two_antenna_port, pmi_typeI_single_panel, pmi_typeII>;
\end{lstlisting}

\subsection{Das komplette Codebook für die Zelle}

\file{lib/fapi\_adaptor/precoding\_codebook\_generator.cpp}, Funktion \code{generate\_precoding\_codebooks()}.
Beim Start wird für die Zelle einmal jede PMI-Kombination durchgezählt. In \code{generate\_pdsch\_sp\_type1()}
laufen vier verschachtelte Schleifen über $i_{1,1}, i_{1,2}, i_{1,3}, i_2$, pro Layer-Anzahl 1 bis
\code{nof\_ports}, und jede Matrix landet unter einem Precoding-Matrix-Index im Repository.

\begin{lstlisting}[language=C++,title={lib/fapi\_adaptor/precoding\_codebook\_generator.cpp, generate\_pdsch\_sp\_type1()}]
pmi_typeI_single_panel_param_ranges param_ranges = get_pmi_ranges_typeI_single_panel(panel, nof_layers);

for (uint8_t i_1_1 = 0; i_1_1 != param_ranges.i_1_1; ++i_1_1) {
  for (uint8_t i_1_2 = 0; i_1_2 != param_ranges.i_1_2; ++i_1_2) {
    for (uint8_t i_1_3 = 0; i_1_3 != param_ranges.i_1_3; ++i_1_3) {
      for (uint8_t i_2 = 0; i_2 != param_ranges.i_2; ++i_2) {
        pmi_typeI_single_panel pmi = {.panel_config = panel,
                                      .i_1_1        = i_1_1,
                                      .i_1_2        = (param_ranges.i_1_2 > 0) ? std::optional(i_1_2) : std::nullopt,
                                      .i_1_3        = (param_ranges.i_1_3 > 0) ? std::optional(i_1_3) : std::nullopt,
                                      .i_2          = i_2};

        unsigned pm_index = base_offset + get_pdsch_single_panel_type1_precoding_matrix_index(param_ranges, pmi);
        repo_builder.add(pm_index, make_composite(pmi, nof_layers, topology));
      }
    }
  }
}
\end{lstlisting}

Daneben werden Identität, SSB, PDCCH, omnidirektionales PDSCH, CSI-RS und PRS eingetragen.
\code{make\_composite()} entscheidet anhand der Antennentopologie, ob die volle Port-Matrix oder die
Beam-Zerlegung gespeichert wird:

\begin{lstlisting}[language=C++,title={lib/fapi\_adaptor/precoding\_codebook\_generator.cpp, make\_composite()}]
static precoding_beamforming_composite
make_composite(const precoding_matrix_indicator& pmi, unsigned nof_layers, antenna_topology topology)
{
  if (has_beam_grid(topology)) {
    return get_mimo_matrix_from_pmi(pmi, nof_layers);
  }

  precoding_weight_matrix precoding = make_precoding(pmi, nof_layers);

  return {precoding, get_default_beam_list(precoding.get_nof_ports())};
}
\end{lstlisting}

Aufgerufen wird die Generierung in \file{lib/du/du\_high/o\_du\_high\_factory.cpp:105} (MAC-Seite, Mapper) und
\file{lib/du/du\_low/o\_du\_low\_factory.cpp:64} (PHY-Seite, Repository). Welches Codebook gilt, hängt nur von
der Anzahl der Sendeantennen ab:

\begin{lstlisting}[language=C++,title={lib/du/du\_low/o\_du\_low\_factory.cpp}]
switch (fapi_cfg.sectors[i].p7_config.carrier_cfg.num_tx_ant) {
  case 1: codebook_config = pmi_codebook_one_port{}; break;
  case 2: codebook_config = pmi_codebook_two_port{}; break;
  case 4:
    codebook_config =
        pmi_codebook_typeI_single_panel{pmi_codebook_single_panel_config::two_one, pmi_codebook_typeI_mode::one};
    break;
  case 8:
    codebook_config =
        pmi_codebook_typeI_single_panel{pmi_codebook_single_panel_config::four_one, pmi_codebook_typeI_mode::one};
    break;
}
\end{lstlisting}

\subsection{Konfiguration und Nutzung im Scheduler}

\begin{itemize}
  \item \file{lib/scheduler/config/csi\_helper.cpp}, \code{make\_type1\_codebook\_config()}: baut die
    RRC-CodebookConfig für das UE (Subset-Restriction-Bitmaps, Panel $2\times1$ oder $4\times1$).
  \item \file{lib/scheduler/ue\_context/ue\_channel\_state\_manager.cpp}: Default-PMI-Konfigurationen und
    Speicherung des vom UE gemeldeten PMI.
\end{itemize}

Zum Einstieg empfiehlt sich die Reihenfolge \code{make\_type1\_sp\_mode1()} $\rightarrow$
\code{generate\_pdsch\_sp\_type1()} $\rightarrow$ \code{calculate\_mimo\_matrix()}.
